\chapter{前言}

我一直在想，這本《人妻約會指南》要寫成一本什麼樣的書？

如果寫成小說，怎麼能超過歌德的《少年維特的煩惱》，又如何能超過史考特·費茲傑羅的《大亨小傳》呢？兩百多年前拿破崙·波拿巴把前者讀了七遍。而後者是現在每個美國高中生的必讀小說。

作為NOI\footnote{中國全國信息學（電腦）奧林匹克競賽}全國第二名和一個頂級程式設計師，我能把這本書寫成像C語言代碼一樣一個分號都不出錯的技術文檔或者操作手冊——但是那樣子，未免也太枯燥一點了！怎麼說，在這本教大家如何和人妻交往的手冊裏面，也要加入一點故事情節和個人情感。

而且，雖然這是一本書，但是大部分人看到的時候，它也是一個PDF。而作為一個PDF，自然要加上精彩的微信截圖和背德照片了。

2021年，我判斷中國房地產要崩盤，於是把我在北京星河灣三千多萬的房子掛出去出售，然後去深圳找做外貿的朋友，看看怎麼把人民幣匯出去投資美國科技股。陰差陽錯之下，又聯繫上了我已婚的初戀女友郭菊陽。我們高中在一起的時候，我爸在監獄裏面。然後她爸跟她說不要談戀愛，要嫁給有錢人，於是就把我甩了。

在深圳歡樂海岸的星巴克甄選裏面，我望着她的眼睛，發誓一定要奪回我被奪走的一切。幾個月後，郭菊陽跪在白石洲的出租屋裏，求我射在她裏面。

所以具體的故事細節，會作為案例、故事細節和個人情感，讓這本全球頂級AI工程師寫的操作手冊，也有那麼一點文學的味道。

\begin{figure}[H]
\centering
\includegraphics[width=1.8in]{figures/guo_1.jpg}
\caption{我在高中給郭菊陽拍的影片的截圖}
\end{figure}

\begin{figure}[H]
\centering
\includegraphics[width=2.5in]{figures/guo_2.jpeg}
\caption{郭菊陽在網誌發的畢業證照片}
\end{figure}

\begin{figure}[H]
\centering
\includegraphics[width=4in]{figures/guo_3.jpg}
\caption{郭菊陽發微博炫耀她找的法拉利富二代}
\end{figure}

\begin{figure}[H]
\centering
\includegraphics[width=3.8in]{figures/guo_4.jpeg}
\caption{郭菊陽知道我更有錢之後，把老公馬賽克了}
\end{figure}

\begin{figure}[H]
\centering
\includegraphics[width=2in]{figures/guo_5.jpeg}
\caption{郭菊陽變成深圳縣城瑜伽褲}
\end{figure}

\begin{figure}[H]
\centering
\includegraphics[width=2in]{figures/guo_6.jpeg}
\caption{郭菊陽給我發的自拍中能過審的一張}
\end{figure}
